\subsection{فعالیت پریدن}

\providecommand{\ResultFigure}[4][0.86\linewidth]{%
\begin{figure}[H]
	\centering
	\includegraphics[width=#1]{#2}
	\caption{#3}
	\label{#4}
\end{figure}}

\subsubsection{ماهیت فعالیت و اهمیت آن}

فعالیت پریدن نسبت به ایستادن و بلند شدن از صندلی، ماهیت سریع‌تر و دینامیک‌تری دارد. این فعالیت شامل چند مرحله اصلی است: آماده‌سازی، اعمال نیرو برای جدا شدن از زمین، فاز پرش، فرود و بازیابی تعادل. به همین دلیل، انتظار می‌رود ویژگی‌هایی مانند دامنه تغییرات سیگنال، انرژی، نوسانات سریع و شاخص‌های مربوط به تغییرات ناگهانی نقش بیشتری در تفکیک سطوح عملکردی داشته باشند.

برای فعالیت پریدن نیز دو حالت داده‌های اصلی و داده‌های جداسازی‌شده با تشخیص فعالیت بررسی شده‌اند. از آنجا که پریدن یک حرکت سریع است، وجود بخش‌های غیرمرتبط قبل و بعد از حرکت می‌تواند اثر بیشتری بر ویژگی‌های آماری بگذارد. بنابراین، مقایسه این دو حالت برای این فعالیت اهمیت ویژه‌ای دارد.

\subsubsection{داده‌های اصلی: سیگنال، کلاس‌ها و ویژگی‌ها}

شکل \ref{fig:jump_normal_signal_sample} نمونه سیگنال‌های ثبت‌شده در فعالیت پریدن را برای داده‌های اصلی نشان می‌دهد. در این فعالیت، نسبت به ایستادن انتظار می‌رود تغییرات سیگنال واضح‌تر و با دامنه بزرگ‌تر باشد. این تغییرات می‌توانند به فاز جهش و فرود مربوط باشند و در ویژگی‌هایی مانند دامنه، انرژی و تغییرات مشتق‌شده اثر بگذارند.

\ResultFigure{Images/Generated/TaskResults/Jump/Normal/signal_sample.png}{نمونه سیگنال‌های ثبت‌شده در فعالیت پریدن، داده‌های اصلی.}{fig:jump_normal_signal_sample}

شکل \ref{fig:jump_normal_class_distribution} توزیع نمونه‌ها را در سطوح عملکردی مختلف نشان می‌دهد. این شکل باید هنگام تفسیر معیارهای کلی مدل در نظر گرفته شود، زیرا عملکرد خوب مدل باید در همه کلاس‌ها بررسی شود.

\ResultFigure[0.72\linewidth]{Images/Generated/TaskResults/Jump/Normal/class_distribution.png}{توزیع نمونه‌ها در فعالیت پریدن، داده‌های اصلی.}{fig:jump_normal_class_distribution}

نقشه همبستگی ویژگی‌ها در شکل \ref{fig:jump_normal_feature_correlation} نشان می‌دهد که بخشی از ویژگی‌های استخراج‌شده دارای وابستگی بالا هستند. در فعالیت پریدن، این همبستگی ممکن است به دلیل ثبت هم‌زمان رویدادهای حرکتی شدید در چند محور یا چند حسگر ایجاد شود.

\ResultFigure{Images/Generated/TaskResults/Jump/Normal/feature_correlation_heatmap.png}{نقشه همبستگی ویژگی‌ها در فعالیت پریدن، داده‌های اصلی.}{fig:jump_normal_feature_correlation}

\subsubsection{داده‌های اصلی: تحلیل \lr{PCA} و مدل‌ها}

شکل \ref{fig:jump_normal_pca_variance} سهم مؤلفه‌های اصلی در توضیح واریانس داده را نشان می‌دهد. این نمودار مشخص می‌کند که چه مقدار از اطلاعات داده با تعداد محدودی مؤلفه قابل نگهداری است. شکل \ref{fig:jump_normal_pca_scatter} نیز پراکندگی نمونه‌ها را در فضای دو مؤلفه اول نشان می‌دهد و دید اولیه‌ای از جدایی یا هم‌پوشانی کلاس‌ها ارائه می‌کند.

\ResultFigure{Images/Generated/TaskResults/Jump/Normal/pca_variance_curve.png}{نمودار واریانس توضیح‌داده‌شده \lr{PCA} در فعالیت پریدن، داده‌های اصلی.}{fig:jump_normal_pca_variance}

\ResultFigure[0.74\linewidth]{Images/Generated/TaskResults/Jump/Normal/pca_scatter_pc1_pc2.png}{پراکندگی نمونه‌های فعالیت پریدن در فضای دو مؤلفه اصلی اول، داده‌های اصلی.}{fig:jump_normal_pca_scatter}

شکل‌های \ref{fig:jump_normal_model_comparison_f1} تا \ref{fig:jump_normal_metric_boxplots} نتایج مدل‌ها را در داده‌های اصلی نمایش می‌دهند. در تفسیر این شکل‌ها باید توجه شود که پریدن یک حرکت چندفازی است و ممکن است مدل‌ها نسبت به ویژگی‌هایی که فازهای مختلف حرکت را به طور جداگانه توصیف می‌کنند حساس باشند.

\ResultFigure{Images/Generated/TaskResults/Jump/Normal/model_comparison_f1.png}{مقایسه عملکرد مدل‌ها بر اساس \lr{F1} در فعالیت پریدن، داده‌های اصلی.}{fig:jump_normal_model_comparison_f1}

\ResultFigure[0.72\linewidth]{Images/Generated/TaskResults/Jump/Normal/model_ranking_f1.png}{رتبه‌بندی مدل‌ها در فعالیت پریدن، داده‌های اصلی.}{fig:jump_normal_model_ranking_f1}

\ResultFigure[0.72\linewidth]{Images/Generated/TaskResults/Jump/Normal/f1_gain_pca_vs_no_pca.png}{اثر \lr{PCA95} بر معیار \lr{F1} در فعالیت پریدن، داده‌های اصلی.}{fig:jump_normal_f1_gain_pca}

\ResultFigure{Images/Generated/TaskResults/Jump/Normal/metric_boxplots.png}{پراکندگی معیارهای ارزیابی در foldهای مختلف، داده‌های اصلی.}{fig:jump_normal_metric_boxplots}

\subsubsection{داده‌های اصلی: تحلیل خطا و ویژگی‌ها}

ماتریس‌های درهم‌ریختگی در شکل‌های \ref{fig:jump_normal_confusion_no_pca} و \ref{fig:jump_normal_confusion_pca95} نشان می‌دهند که خطاهای مدل در کدام کلاس‌ها متمرکز بوده‌اند. در فعالیت پریدن، اگر خطاها بین کلاس‌های مجاور رخ دهند، می‌تواند به شباهت نسبی الگوی حرکتی سطوح عملکردی نزدیک مربوط باشد. اما اگر خطاهای غیرمجاور دیده شود، لازم است کیفیت داده یا ویژگی‌ها با دقت بیشتری بررسی شود.

\ResultFigure{Images/Generated/TaskResults/Jump/Normal/confusion_matrices_no_pca.png}{ماتریس‌های درهم‌ریختگی در فعالیت پریدن، داده‌های اصلی، بدون \lr{PCA}.}{fig:jump_normal_confusion_no_pca}

\ResultFigure{Images/Generated/TaskResults/Jump/Normal/per_class_recall_no_pca.png}{مقایسه \lr{Recall} هر کلاس در فعالیت پریدن، داده‌های اصلی، بدون \lr{PCA}.}{fig:jump_normal_per_class_recall_no_pca}

\ResultFigure{Images/Generated/TaskResults/Jump/Normal/confusion_matrices_pca95.png}{ماتریس‌های درهم‌ریختگی در فعالیت پریدن، داده‌های اصلی، با \lr{PCA95}.}{fig:jump_normal_confusion_pca95}

\ResultFigure{Images/Generated/TaskResults/Jump/Normal/per_class_recall_pca95.png}{مقایسه \lr{Recall} هر کلاس در فعالیت پریدن، داده‌های اصلی، با \lr{PCA95}.}{fig:jump_normal_per_class_recall_pca95}

شکل \ref{fig:jump_normal_rf_importance} اهمیت ویژگی‌ها را بر اساس مدل \lr{Random Forest} نشان می‌دهد. در فعالیت پریدن، انتظار می‌رود ویژگی‌های مرتبط با پاها، دامنه تغییرات و انرژی سیگنال نقش پررنگی داشته باشند. اگر ویژگی‌های سر یا دست‌ها نیز در میان ویژگی‌های مهم دیده شوند، می‌توان آن را به نقش تعادل، وضعیت تنه و حرکات جبرانی نسبت داد.

\ResultFigure{Images/Generated/TaskResults/Jump/Normal/rf_top_feature_importance.png}{ویژگی‌های مهم‌تر بر اساس \lr{Random Forest}، داده‌های اصلی.}{fig:jump_normal_rf_importance}

\ResultFigure{Images/Generated/TaskResults/Jump/Normal/class_centroid_feature_heatmap.png}{میانگین نرمال‌شده ویژگی‌های منتخب در کلاس‌ها، داده‌های اصلی.}{fig:jump_normal_class_centroid_heatmap}

\ResultFigure{Images/Generated/TaskResults/Jump/Normal/svm_metrics_vs_pca_components.png}{تغییر عملکرد \lr{SVM} نسبت به تعداد مؤلفه‌های اصلی، داده‌های اصلی.}{fig:jump_normal_svm_metrics_vs_pca_components}

\ResultFigure{Images/Generated/TaskResults/Jump/Normal/models_f1_vs_pca_components.png}{مقایسه مدل‌ها نسبت به تعداد مؤلفه‌های اصلی، داده‌های اصلی.}{fig:jump_normal_models_f1_vs_pca_components}

\subsubsection{داده‌های جداسازی‌شده با تشخیص فعالیت}

در داده‌های جداسازی‌شده، هدف تمرکز ویژگی‌ها بر بازه‌ای است که واقعاً حرکت پریدن در آن رخ داده است. شکل \ref{fig:jump_detectedaction_signal_sample} نمونه سیگنال این حالت را نشان می‌دهد. از آنجا که پریدن سریع است، جداسازی دقیق بازه حرکت می‌تواند اثر زیادی بر ویژگی‌های آماری داشته باشد.

\ResultFigure{Images/Generated/TaskResults/Jump/DetectedAction/signal_sample.png}{نمونه سیگنال‌های ثبت‌شده در فعالیت پریدن، داده‌های جداسازی‌شده.}{fig:jump_detectedaction_signal_sample}

\ResultFigure[0.72\linewidth]{Images/Generated/TaskResults/Jump/DetectedAction/class_distribution.png}{توزیع نمونه‌ها در فعالیت پریدن، داده‌های جداسازی‌شده.}{fig:jump_detectedaction_class_distribution}

\ResultFigure{Images/Generated/TaskResults/Jump/DetectedAction/feature_correlation_heatmap.png}{نقشه همبستگی ویژگی‌ها در فعالیت پریدن، داده‌های جداسازی‌شده.}{fig:jump_detectedaction_feature_correlation}

\ResultFigure{Images/Generated/TaskResults/Jump/DetectedAction/pca_variance_curve.png}{نمودار واریانس توضیح‌داده‌شده \lr{PCA}، داده‌های جداسازی‌شده.}{fig:jump_detectedaction_pca_variance}

\ResultFigure[0.74\linewidth]{Images/Generated/TaskResults/Jump/DetectedAction/pca_scatter_pc1_pc2.png}{پراکندگی نمونه‌ها در فضای دو مؤلفه اصلی اول، داده‌های جداسازی‌شده.}{fig:jump_detectedaction_pca_scatter}

شکل‌های زیر عملکرد مدل‌ها را پس از جداسازی بازه فعالیت نشان می‌دهند. مقایسه این خروجی‌ها با داده‌های اصلی می‌تواند مشخص کند که آیا حذف بخش‌های غیرمرتبط باعث بهبود عملکرد مدل شده است یا خیر.

\ResultFigure{Images/Generated/TaskResults/Jump/DetectedAction/model_comparison_f1.png}{مقایسه عملکرد مدل‌ها بر اساس \lr{F1}، داده‌های جداسازی‌شده.}{fig:jump_detectedaction_model_comparison_f1}

\ResultFigure[0.72\linewidth]{Images/Generated/TaskResults/Jump/DetectedAction/model_ranking_f1.png}{رتبه‌بندی مدل‌ها، داده‌های جداسازی‌شده.}{fig:jump_detectedaction_model_ranking_f1}

\ResultFigure[0.72\linewidth]{Images/Generated/TaskResults/Jump/DetectedAction/f1_gain_pca_vs_no_pca.png}{اثر \lr{PCA95} بر معیار \lr{F1}، داده‌های جداسازی‌شده.}{fig:jump_detectedaction_f1_gain_pca}

\ResultFigure{Images/Generated/TaskResults/Jump/DetectedAction/metric_boxplots.png}{پراکندگی معیارهای ارزیابی در foldهای مختلف، داده‌های جداسازی‌شده.}{fig:jump_detectedaction_metric_boxplots}

\ResultFigure{Images/Generated/TaskResults/Jump/DetectedAction/confusion_matrices_no_pca.png}{ماتریس‌های درهم‌ریختگی، داده‌های جداسازی‌شده، بدون \lr{PCA}.}{fig:jump_detectedaction_confusion_no_pca}

\ResultFigure{Images/Generated/TaskResults/Jump/DetectedAction/per_class_recall_no_pca.png}{مقایسه \lr{Recall} هر کلاس، داده‌های جداسازی‌شده، بدون \lr{PCA}.}{fig:jump_detectedaction_per_class_recall_no_pca}

\ResultFigure{Images/Generated/TaskResults/Jump/DetectedAction/confusion_matrices_pca95.png}{ماتریس‌های درهم‌ریختگی، داده‌های جداسازی‌شده، با \lr{PCA95}.}{fig:jump_detectedaction_confusion_pca95}

\ResultFigure{Images/Generated/TaskResults/Jump/DetectedAction/per_class_recall_pca95.png}{مقایسه \lr{Recall} هر کلاس، داده‌های جداسازی‌شده، با \lr{PCA95}.}{fig:jump_detectedaction_per_class_recall_pca95}

\ResultFigure{Images/Generated/TaskResults/Jump/DetectedAction/rf_top_feature_importance.png}{ویژگی‌های مهم‌تر بر اساس \lr{Random Forest}، داده‌های جداسازی‌شده.}{fig:jump_detectedaction_rf_importance}

\ResultFigure{Images/Generated/TaskResults/Jump/DetectedAction/class_centroid_feature_heatmap.png}{میانگین نرمال‌شده ویژگی‌های منتخب در کلاس‌ها، داده‌های جداسازی‌شده.}{fig:jump_detectedaction_class_centroid_heatmap}

\ResultFigure{Images/Generated/TaskResults/Jump/DetectedAction/svm_metrics_vs_pca_components.png}{تغییر عملکرد \lr{SVM} نسبت به تعداد مؤلفه‌های اصلی، داده‌های جداسازی‌شده.}{fig:jump_detectedaction_svm_metrics_vs_pca_components}

\ResultFigure{Images/Generated/TaskResults/Jump/DetectedAction/models_f1_vs_pca_components.png}{مقایسه مدل‌ها نسبت به تعداد مؤلفه‌های اصلی، داده‌های جداسازی‌شده.}{fig:jump_detectedaction_models_f1_vs_pca_components}

% TODO: عددهای جدول زیر را از خروجی‌های Excel فعالیت Jump وارد کن.
\begin{table}[H]
	\centering
	\caption{خلاصه مقایسه داده‌های اصلی و جداسازی‌شده در فعالیت پریدن}
	\label{tab:jump_summary}
	\begin{tabular}{p{3cm}lccccc}
		\toprule
		نوع داده & مدل & \lr{Accuracy} & \lr{Precision} & \lr{Recall} & \lr{Specificity} & \lr{F1} \\
		\midrule
		داده‌های اصلی & \lr{KNN} &  &  &  &  &  \\
		داده‌های اصلی & \lr{SVM} &  &  &  &  &  \\
		داده‌های اصلی & \lr{Random Forest} &  &  &  &  &  \\
		داده‌های جداسازی‌شده & \lr{KNN} &  &  &  &  &  \\
		داده‌های جداسازی‌شده & \lr{SVM} &  &  &  &  &  \\
		داده‌های جداسازی‌شده & \lr{Random Forest} &  &  &  &  &  \\
		\bottomrule
	\end{tabular}
\end{table}

\subsubsection{جمع‌بندی فعالیت پریدن}

فعالیت پریدن به دلیل ماهیت دینامیک خود، یکی از فعالیت‌های مهم برای ارزیابی توان سامانه در تحلیل حرکات سریع است. در این فعالیت، مقایسه داده‌های اصلی و داده‌های جداسازی‌شده اهمیت زیادی دارد، زیرا بازه‌های قبل و بعد از پرش می‌توانند بر ویژگی‌های آماری اثر بگذارند. پس از تکمیل جدول عددی، باید مشخص شود کدام حالت داده و کدام مدل برای این فعالیت مناسب‌تر است و آیا جداسازی بازه حرکت باعث بهبود عملکرد شده است یا خیر.
